% =====================================================================
%  BAB 3 — METODOLOGI PENELITIAN
% =====================================================================

\chapter{Metodologi Penelitian}

\section{Tipe Penelitian}
\label{sec:3.1}

Penelitian ini merupakan penelitian non-implementatif analitik dengan
pendekatan kuantitatif, karena berfokus pada analisis pengaruh strategi
\emph{fine-tuning} terhadap akurasi klasifikasi emosi ucapan, bukan pada
pengembangan sistem atau perangkat lunak siap pakai. Sifat analitik
ditunjukkan melalui pengolahan data hasil observasi akurasi dan F1-score
dari setiap strategi \emph{fine-tuning} untuk mengkaji hubungan antara
strategi yang diterapkan dan performa klasifikasi emosi yang dihasilkan.
Pendekatan kuantitatif digunakan karena data berupa nilai numerik hasil
evaluasi model yang dianalisis menggunakan metrik akurasi, F1-score, dan
pengujian signifikansi statistik. Model-model yang di-\emph{fine-tune}
pada penelitian ini berfungsi sebagai instrumen untuk menghasilkan data
observasi tersebut, bukan sebagai produk/artefak utama penelitian.

\begin{figure}[H]
  \centering
  \begin{tikzpicture}[node distance=8mm]
    \node[term] (mulai)   {Mulai};
    \node[proc] (studi)   [below=of mulai]    {Studi Literatur};
    \node[proc] (data)    [below=of studi]    {Pengumpulan Data};
    \node[proc] (sumber)  [below=of data]     {Fine-Tuning Sumber};
    \node[proc] (target)  [below=of sumber]   {Fine-Tuning Target};
    \node[proc] (evaluasi)[below=of target]   {Evaluasi K-Fold};
    \node[proc] (kinerja) [below=of evaluasi] {Analisis Kinerja Komputasi};
    \node[proc] (analisis)[below=of kinerja]  {Analisis Hasil};
    \node[term] (selesai) [below=of analisis] {Selesai};

    \draw[arr] (mulai)    -- (studi);
    \draw[arr] (studi)    -- (data);
    \draw[arr] (data)     -- (sumber);
    \draw[arr] (sumber)   -- (target);
    \draw[arr] (target)   -- (evaluasi);
    \draw[arr] (evaluasi) -- (kinerja);
    \draw[arr] (kinerja)  -- (analisis);
    \draw[arr] (analisis) -- (selesai);
  \end{tikzpicture}
  \caption{Diagram alir metode penelitian}
  \label{fig:3.alur}
\end{figure}

Diagram alir metode penelitian ini ditunjukkan pada Gambar
\ref{fig:3.alur}. Tahap awal penelitian diawali dengan studi literatur,
yang bertujuan untuk mengkaji teori, metode, serta penelitian terdahulu
yang relevan dengan SER, pembelajaran \emph{self-supervised}, dan
\emph{cross-lingual transfer}. Tahap berikutnya adalah pengumpulan data,
yaitu proses pengumpulan korpus emosi berbahasa Inggris sebagai sumber
dan dataset SER berbahasa Indonesia sebagai target. Data yang telah
diperoleh selanjutnya diproses pada tahap \emph{fine-tuning} sumber
untuk melatih backbone \emph{self-supervised} mengenali pola akustik
emosi secara umum. Setelah itu dilakukan \emph{fine-tuning} target untuk
mengadaptasi model dasar ke data Indonesia melalui beberapa skema
strategi. Tahap berikutnya adalah evaluasi k-fold untuk menilai performa
setiap strategi menggunakan skema validasi silang dan pengujian
signifikansi statistik, dilanjutkan analisis kinerja komputasi untuk
membandingkan waktu pelatihan dan penggunaan memori antarkondisi. Hasil
kedua evaluasi tersebut kemudian digunakan sebagai dasar pada tahap
analisis hasil, yang merangkum temuan penelitian serta menjawab rumusan
masalah yang telah dipaparkan.

\section{Studi Literatur}
\label{sec:3.2}

Studi literatur merupakan tahap awal penelitian yang bertujuan untuk
mempelajari dan memahami teori, konsep, serta hasil penelitian terdahulu
yang berkaitan dengan SER, pembelajaran \emph{self-supervised},
representasi Wav2Vec2/XLSR, dan \emph{cross-lingual transfer}. Pada
tahap ini dilakukan pengkajian terhadap jurnal, prosiding konferensi,
dan sumber ilmiah relevan untuk memastikan celah penelitian pada Bab 2
Landasan Kepustakaan belum terjawab oleh penelitian lain, sekaligus
menentukan backbone model dan skema \emph{fine-tuning} yang sesuai
dengan tujuan penelitian.

\section{Pengumpulan Data}
\label{sec:3.3}

Tahap ini mengumpulkan dua kelompok data, yaitu korpus emosi berbahasa
Inggris sebagai sumber pada tahap \emph{fine-tuning} pertama
(RAVDESS/CREMA-D), dan dataset E-SERAVD (Santoso, Budianto and Dutono,
2026) sebagai target berbahasa Indonesia pada tahap \emph{fine-tuning}
kedua. Skema kelas emosi pada korpus sumber disesuaikan dengan skema
empat kelas yang tersedia pada E-SERAVD (marah, senang, netral, sedih),
sehingga perbandingan antarstrategi \emph{fine-tuning} dapat dilakukan
secara adil.

Karakteristik masing-masing korpus dijelaskan sebagai berikut. RAVDESS
adalah korpus ucapan emosional berbahasa Inggris yang direkam oleh 24
aktor profesional (12 perempuan, 12 laki-laki), mencakup 8 kategori
emosi (netral, tenang, senang, sedih, marah, takut, jijik, dan terkejut)
(Livingstone and Russo, 2018). CREMA-D adalah korpus ucapan emosional
berbahasa Inggris yang direkam oleh 91 aktor dengan latar belakang usia
dan etnis yang beragam, mencakup 6 kategori emosi (marah, jijik, takut,
senang, netral, dan sedih) (Cao et al., 2014). Kedua korpus ini dipilih
sebagai sumber karena tersedia secara publik, telah banyak digunakan
pada penelitian SER sebelumnya, dan mencakup kategori emosi yang secara
umum beririsan dengan skema kelas pada dataset target. Tabel
\ref{tab:3.dataset} merangkum karakteristik ketiga dataset yang
digunakan.

\begin{table}[H]
  \centering
  \caption{Karakteristik dataset yang digunakan}
  \label{tab:3.dataset}
  \begin{adjustbox}{max width=\textwidth}
  \begin{tabular}{p{3.8cm} l l p{2.6cm} p{1.8cm}}
    \toprule
    \textbf{Dataset} & \textbf{Peran} & \textbf{Bahasa} & \textbf{Jumlah Sampel} & \textbf{Jumlah Kelas} \\
    \midrule
    RAVDESS & Sumber & Inggris & 1.440 & 8 \\
    CREMA-D & Sumber & Inggris & 7.442 & 6 \\
    E-SERAVD & Target & Indonesia & $\sim$5.049 & 4 \\
    \bottomrule
  \end{tabular}
  \end{adjustbox}
\end{table}

\noindent
Jumlah sampel untuk RAVDESS dan CREMA-D di atas mengacu pada dokumentasi
resmi masing-masing korpus. E-SERAVD memiliki sekitar 5.049 sampel
dengan distribusi 1.202 marah, 1.228 senang, 2.075 netral, dan 899 sedih
(Santoso, Budianto and Dutono, 2026); angka ini masih perlu diverifikasi
ulang terhadap dokumentasi resmi dataset sebelum benar-benar diperoleh.
Dataset SER Indonesia lain seperti IndoWaveSentiment (Nugroho et al.,
2025) tidak digunakan sebagai target pada penelitian ini karena skala
dan skema kelasnya berbeda dari E-SERAVD, namun disebutkan kembali
sebagai arah pengembangan pada Bab 6.

\section{Fine-Tuning Sumber}
\label{sec:3.4}

Backbone yang digunakan adalah model dasar XLSR-53
(\path{facebook/wav2vec2-large-xlsr-53}) tanpa \emph{fine-tuning}
tambahan dari pihak lain untuk Bahasa Indonesia -- berbeda dari
pendekatan Nugroho et al. (2025) yang menguji model dasar XLSR-Wav2Vec2
yang sudah di-\emph{fine-tune} untuk ucapan Indonesia (lihat Subbab
\ref{sec:2.4}). Pemilihan ini disengaja: dengan titik awal yang identik
di keempat kondisi eksperimen (Subbab \ref{sec:3.5}), perbedaan performa
yang teramati dapat ditarik ke strategi \emph{fine-tuning} yang diuji
dalam penelitian ini, bukan tercampur dengan efek \emph{fine-tuning}
pihak lain yang sudah menempel pada checkpoint.

Tahap ini melatih backbone \emph{self-supervised} (Wav2Vec2/XLSR-53)
pada korpus emosi berbahasa Inggris agar model mampu mengenali pola
akustik emosi secara umum, menghasilkan model dasar yang selanjutnya
akan diadaptasi ke data target pada tahap berikutnya. Tahap ini
dijalankan sekali saja; bobot hasilnya dipakai ulang sebagai titik awal
untuk semua fold pada Kondisi C dan D (Subbab \ref{sec:3.5}), bukan
diulang per-fold, karena \emph{fine-tuning} tahap sumber tidak
bergantung pada pembagian fold data target.

\section{Fine-Tuning Target}
\label{sec:3.5}

Tahap ini mengadaptasi model ke data Indonesia melalui dua axis desain
yang independen: (1) \emph{stage}, yaitu apakah model melalui
\emph{fine-tuning} dua tahap atau langsung dilatih pada data target
(\emph{single-stage}); dan (2) metode adaptasi, yaitu \emph{full
fine-tuning} yang memperbarui seluruh parameter model, atau LoRA yang
hanya memperbarui parameter tambahan berjumlah kecil. Kombinasi kedua
axis tersebut menghasilkan empat kondisi eksperimen sebagaimana
dirangkum pada Tabel \ref{tab:3.kondisi}.

\begin{table}[H]
  \centering
  \caption{Empat kondisi eksperimen (desain 2$\times$2)}
  \label{tab:3.kondisi}
  \begin{tabular}{c l l}
    \toprule
    \textbf{Kondisi} & \textbf{Stage} & \textbf{Metode Adaptasi} \\
    \midrule
    A & Single-stage & Full Fine-Tuning \\
    B & Single-stage & LoRA \\
    C & Two-stage    & Full Fine-Tuning \\
    D & Two-stage    & LoRA \\
    \bottomrule
  \end{tabular}
\end{table}

\noindent
Backbone yang digunakan sama persis (XLSR-53, Subbab \ref{sec:3.4}) di
keempat kondisi -- yang berbeda hanya resep \emph{fine-tuning}-nya --
supaya perbedaan hasil antarkondisi dapat ditarik ke strategi
\emph{fine-tuning} yang diuji, bukan ke perbedaan arsitektur atau titik
awal model. Kondisi A dan B tidak melalui tahap \emph{fine-tuning}
sumber (Subbab \ref{sec:3.4}) dan langsung dilatih pada data target,
sedangkan Kondisi C dan D memakai bobot hasil \emph{fine-tuning} sumber
sebagai titik awal. Validasi silang k-fold (Subbab \ref{sec:3.6})
diterapkan pada level adaptasi target ini untuk keempat kondisi. Sebagai
mitigasi risiko, direkomendasikan menjalankan \emph{pilot run} (1
kondisi $\times$ 1 fold) pada perangkat keras sebenarnya sebelum
menjalankan grid eksperimen penuh, untuk memverifikasi estimasi waktu
dan VRAM pada Subbab \ref{sec:3.10.2} dengan data riil.

\section{Evaluasi K-Fold}
\label{sec:3.6}

Tahap ini mengevaluasi setiap kondisi eksperimen menggunakan skema
validasi silang k-fold, sehingga diperoleh beberapa observasi akurasi
dan F1-score per fold untuk masing-masing kondisi. Data hasil observasi
tersebut kemudian diuji normalitasnya menggunakan uji Shapiro-Wilk. Jika
data berdistribusi normal, pengujian dilanjutkan dengan \emph{paired
t-test}; jika tidak, digunakan \emph{Wilcoxon signed-rank test}.
Pengujian ini bertujuan untuk mengetahui signifikansi statistik
perbedaan performa antarkondisi, menjawab hipotesis yang telah
dirumuskan pada Subbab 1.3.

\section{Analisis Kinerja Komputasi}
\label{sec:3.7}

Tahap ini menganalisis kinerja komputasi setiap kondisi eksperimen
(Kondisi A--D), meliputi waktu pelatihan dan penggunaan memori (VRAM)
sebagaimana didefinisikan pada Subbab \ref{sec:2.9.2}. Waktu pelatihan
dicatat dari awal hingga akhir proses pelatihan pada tiap fold untuk
tahap \emph{fine-tuning} target (Subbab \ref{sec:3.5}); durasi
\emph{fine-tuning} sumber (Subbab \ref{sec:3.4}) dicatat terpisah karena
hanya dijalankan sekali dan dipakai bersama oleh Kondisi C dan D.
Penggunaan memori diukur sebagai puncak alokasi VRAM selama pelatihan
pada tiap kondisi. Secara teknis, waktu pelatihan dicatat menggunakan
\texttt{time.perf\_counter()}, sedangkan puncak penggunaan VRAM dicatat
menggunakan \texttt{torch.cuda.max\_memory\_allocated()} setelah
\texttt{reset\_peak\_memory\_stats()} dipanggil di awal tiap fold. Kedua
metrik ini dibandingkan secara deskriptif
antarkondisi, dengan perhatian khusus pada perbandingan Kondisi A vs B
dan Kondisi C vs D (axis \emph{full fine-tuning} vs LoRA) untuk menjawab
Rumusan Masalah 2 dari sisi efisiensi, melengkapi pengujian signifikansi
statistik akurasi dan F1-score pada Subbab \ref{sec:3.6}.

\section{Analisis Hasil}
\label{sec:3.8}

Tahap ini menafsirkan hasil pengujian signifikansi statistik dan
analisis kinerja komputasi untuk menjawab rumusan masalah penelitian,
serta menyusun rekomendasi strategi \emph{fine-tuning} yang paling
sesuai untuk pengembangan SER berbahasa Indonesia berdasarkan temuan
yang diperoleh.

\section{Kesimpulan dan Saran}
\label{sec:3.9}

Penarikan kesimpulan dilaksanakan setelah seluruh tahap evaluasi dan
analisis hasil selesai. Kesimpulan memuat jawaban atas rumusan masalah
yang telah dituliskan sebelumnya, mengacu pada hasil pengujian
signifikansi statistik dari setiap strategi \emph{fine-tuning} yang
dibandingkan. Selain kesimpulan, disusun pula saran sebagai bahan
pertimbangan dan perbaikan untuk penelitian selanjutnya, misalnya
perluasan cakupan korpus sumber atau backbone model yang diuji.

\section{Peralatan Pendukung}
\label{sec:3.10}

Peralatan pendukung yang digunakan dalam penelitian ini terdiri atas
perangkat lunak (\emph{software}) dan perangkat keras (\emph{hardware}).

\subsection{Perangkat Lunak (Software)}
\label{sec:3.10.1}

Perangkat lunak yang digunakan dalam mendukung penelitian ini sebagai
berikut.

\begin{enumerate}
  \item Bahasa pemrograman Python.
  \item Kerangka kerja \emph{deep learning} untuk pelatihan dan
        \emph{fine-tuning} model Wav2Vec2/XLSR.
  \item Pustaka untuk implementasi LoRA.
  \item Pustaka optimizer 8-bit (\texttt{bitsandbytes}) untuk mitigasi
        VRAM pada kondisi \emph{full fine-tuning}.
  \item Pustaka statistik untuk pengujian normalitas dan uji
        signifikansi.
  \item GitHub untuk pengelolaan versi kode sumber dan eksperimen.
\end{enumerate}

\subsection{Perangkat Keras (Hardware)}
\label{sec:3.10.2}

Perangkat keras yang digunakan dalam mendukung penelitian ini sebagai
berikut.

\begin{enumerate}
  \item 2 unit GPU NVIDIA GeForce RTX 3060 6GB (server lab, dipinjamkan).
  \item RAM dan penyimpanan: [\emph{diisi sesuai spesifikasi server lab
        yang sebenarnya dipakai}].
\end{enumerate}

Kecukupan VRAM \emph{per kartu} dianalisis terhadap ukuran backbone
(sekitar 317 juta parameter), sebagaimana dirangkum pada Tabel
\ref{tab:3.vram}.

\begin{table}[H]
  \centering
  \caption{Estimasi kecukupan VRAM per kartu (6GB)}
  \label{tab:3.vram}
  \begin{adjustbox}{max width=\textwidth}
  \begin{tabular}{l l l l}
    \toprule
    \textbf{Komponen} & \textbf{Full-FT (naif)} & \textbf{Full-FT + 8-bit Adam} & \textbf{LoRA} \\
    \midrule
    Bobot + master + gradien + optimizer & $\sim$5 GB & $\sim$1,9 GB & $\sim$0,6--0,7 GB \\
    Sisa untuk aktivasi (dari 6 GB)      & $\sim$1 GB & $\sim$4 GB   & $\sim$5,3 GB \\
    \bottomrule
  \end{tabular}
  \end{adjustbox}
\end{table}

\noindent
Estimasi ``naif'' (optimizer Adam presisi-penuh standar) menyisakan
hanya sekitar 1 GB untuk aktivasi -- terlalu mepet untuk aman, karena
\emph{overhead} CUDA dan fragmentasi memori berisiko menghabiskan sisa
tersebut sebelum sempat dipakai aktivasi. Mitigasi yang digunakan untuk
Kondisi A dan C (\emph{full fine-tuning}) adalah optimizer Adam 8-bit
(\texttt{bitsandbytes}), yang memangkas kebutuhan memori optimizer
menjadi sekitar 1,9 GB dan menyisakan sekitar 4 GB untuk aktivasi --
jauh lebih aman. Dengan mitigasi ini, \emph{full fine-tuning}
diperkirakan tetap memerlukan \emph{batch size} kecil (sekitar 1--2),
dengan \emph{gradient checkpointing} dan \emph{gradient accumulation}
sebagai cadangan opsional jika 4 GB masih belum cukup pada praktiknya.
Skema LoRA (Kondisi B dan D) memiliki ruang gerak yang jauh lebih
longgar, dengan \emph{batch size} 8--16 yang tetap realistis tanpa
teknik tambahan. Keterbatasan ini sekaligus menjadi ilustrasi konkret
argumen efisiensi LoRA yang diuji dalam penelitian ini (Subbab
\ref{sec:2.7}): pada perangkat keras yang benar-benar terbatas, LoRA
bukan hanya lebih hemat secara teoretis, tetapi menjadi opsi yang jauh
lebih praktis dijalankan.

Strategi 2 GPU yang digunakan bukan \emph{distributed training} (model
sudah muat pada 1 kartu), melainkan 2 \emph{run} independen paralel,
satu per GPU (diatur lewat \texttt{CUDA\_VISIBLE\_DEVICES}), karena grid
eksperimen (4 kondisi $\times$ k fold) bersifat \emph{embarrassingly
parallel}.

Rekomendasi teknis lain: \emph{mixed precision} (fp16/bf16) untuk
mengurangi jejak memori, serta membekukan \emph{CNN feature extractor}
(praktik standar pada \emph{fine-tuning} wav2vec2). Dengan skala dataset
target diperkirakan sekitar 5 ribu sampel (Subbab \ref{sec:3.3}),
estimasi kasar waktu per-\emph{run} berkisar puluhan menit hingga
beberapa jam -- kemungkinan lebih lama untuk kondisi \emph{full
fine-tuning} karena \emph{batch size} yang lebih kecil -- sehingga total
waktu penelitian ($k=5$, sekitar 21 \emph{run}) diperkirakan selesai
dalam hitungan hari dengan pemanfaatan 2 GPU secara paralel. Estimasi
ini bersifat kasar dan perlu diverifikasi pada uji coba awal.
